\documentclass{article}
\usepackage{amsmath,amssymb,amsthm}
\usepackage{algorithm}
\usepackage{algpseudocode}
\usepackage{hyperref}
\usepackage{enumitem}
\newtheorem{theorem}{Theorem}
\newtheorem{lemma}{Lemma}
\newtheorem{assumption}{Assumption}
\newtheorem{corollary}{Corollary}
\newtheorem{remark}{Remark}
\newcommand{\EE}{\mathbb{E}}
\newcommand{\RR}{\mathbb{R}}
\newcommand{\norm}[1]{\left\|#1\right\|}
\newcommand{\inner}[2]{\left\langle #1,#2\right\rangle}
\begin{document}

\section*{Algorithm Name}
\textbf{Snapshot‑Interpolation Stochastic Gradient Descent (SIS‑SGD)}  

The method runs ordinary stochastic gradient descent, stores the iterate every $k$ steps as a \emph{snapshot}, and for any intermediate time uses a linear interpolation of the two neighbouring snapshots as the prediction model.

--------------------------------------------------------------------

\section*{Mathematical Formulation}
Let $f:\RR^{d}\rightarrow\RR$ be the (possibly non‑convex) objective and let
\[
g_t:=\nabla f(w_t;\xi_t)
\]
be a stochastic gradient computed on a random mini‑batch $\xi_t$ at iteration $t$.
\begin{itemize}
    \item \textbf{SGD update}
    \begin{equation}
        w_{t+1}=w_t-\eta_t g_t,\qquad t=0,1,\dots,T-1,
        \label{eq:update}
    \end{equation}
    where $\eta_t>0$ is the learning rate.
    \item \textbf{Snapshot collection}
    \[
        \tilde w_{s}=w_{sk},\qquad s=0,1,\dots,S-1,\quad S:=\Big\lfloor\frac{T}{k}\Big\rfloor .
    \]
    The index $s$ denotes the $s$‑th snapshot.
    \item \textbf{Interpolation for prediction}
    For any $t\in[sk,(s+1)k]$ define
    \begin{equation}
        \alpha_t:=\frac{t-sk}{k}\in[0,1],\qquad 
        \hat w_t:=(1-\alpha_t)\tilde w_{s}+\alpha_t\tilde w_{s+1}.
        \label{eq:interp}
    \end{equation}
    The vector $\hat w_t$ is the model used for inference at time $t$.
\end{itemize}
The algorithm returns either the last interpolated iterate $\hat w_T$ or the simple average of all snapshots $\bar w:=\frac1S\sum_{s=0}^{S-1}\tilde w_s$.

--------------------------------------------------------------------

\section*{Pseudocode}
\begin{algorithm}[H]
\caption{Snapshot‑Interpolation SGD (SIS‑SGD)}
\begin{algorithmic}[1]
\State \textbf{Input:} initial point $w_0$, learning‑rate schedule $\{\eta_t\}$, snapshot period $k$, total iterations $T$.
\State Initialize empty list $\mathcal{S}\leftarrow[\,]$.
\For{$t=0$ \textbf{to} $T-1$}
    \State Sample mini‑batch $\xi_t$ and compute stochastic gradient $g_t=\nabla f(w_t;\xi_t)$.
    \State $w_{t+1}\gets w_t-\eta_t g_t$ \hfill // \eqref{eq:update}
    \If{$(t+1)\bmod k =0$}
        \State Append $w_{t+1}$ to $\mathcal{S}$ as $\tilde w_{s}$  \hfill // snapshot
    \EndIf
\EndFor
\State \textbf{Prediction (any $t$):}
\State $s\gets\lfloor t/k\rfloor$, $\alpha_t\gets (t-sk)/k$,
\State $\hat w_t\gets (1-\alpha_t)\tilde w_{s}+\alpha_t\tilde w_{s+1}$ \hfill // \eqref{eq:interp}
\State \textbf{Output:} $\hat w_T$ (or $\bar w=\frac1{|\mathcal{S}|}\sum_{s}\tilde w_s$).
\end{algorithmic}
\end{algorithm}

--------------------------------------------------------------------

\section*{Dry Run Example}
Consider the scalar quadratic $f(w)=(w-3)^2$, thus $\nabla f(w)=2(w-3)$.  
Set $w_0=0$, constant step size $\eta=0.1$, snapshot period $k=2$, and run $T=3$ iterations.

\begin{center}
\begin{tabular}{c|c|c|c|c}
$t$ & $w_t$ & $g_t=2(w_t-3)$ & $w_{t+1}=w_t-\eta g_t$ & $f(w_{t+1})$\\ \hline
0 & $0$ & $-6$ & $0-0.1(-6)=0.6$ & $(0.6-3)^2=5.76$\\
1 & $0.6$ & $-4.8$ & $0.6-0.1(-4.8)=1.08$ & $(1.08-3)^2=3.6864$\\
2 & $1.08$ & $-3.84$ & $1.08-0.1(-3.84)=1.464$ & $(1.464-3)^2=2.3665$\\
\end{tabular}
\end{center}
Snapshots: $\tilde w_{0}=w_{2}=1.08$ (since $k=2$) and $\tilde w_{1}=w_{3}=1.464$ (the last iterate).  

Interpolation for $t=1$ (between snapshots $0$ and $1$):
\[
\alpha_1=\frac{1-0\cdot2}{2}=0.5,\qquad \hat w_1=0.5\cdot\tilde w_{0}+0.5\cdot\tilde w_{1}=0.5(1.08+1.464)=1.272.
\]
The corresponding function value is $f(\hat w_1)=(1.272-3)^2\approx2.985$.

--------------------------------------------------------------------

\section*{Assumptions}
\begin{assumption}[Smoothness]\label{ass:smooth}
$f$ is $L$‑smooth: for all $x,y\in\RR^{d}$,
\[
f(y)\le f(x)+\inner{\nabla f(x)}{y-x}+\frac{L}{2}\norm{y-x}^2 .
\]
\end{assumption}
\begin{assumption}[Unbiased stochastic gradients]\label{ass:unbiased}
For every $t$, $\EE[g_t\mid w_t]=\nabla f(w_t)$.
\end{assumption}
\begin{assumption}[Bounded variance]\label{ass:var}
There exists $\sigma^2\ge0$ such that
\[
\EE\bigl[\norm{g_t-\nabla f(w_t)}^2\mid w_t\bigr]\le\sigma^2,\qquad\forall t .
\]
\end{assumption}
\begin{assumption}[Learning‑rate]\label{ass:lr}
The step size is constant $\eta_t\equiv\eta$ with $0<\eta\le \frac{1}{L}$.
\end{assumption}
All theorems below rely only on the four assumptions above.

--------------------------------------------------------------------

\section*{Main Convergence Theorem}
\begin{theorem}[Convergence of SIS‑SGD]\label{thm:main}
Under Assumptions \ref{ass:smooth}–\ref{ass:lr}, let $S=\lfloor T/k\rfloor$ be the number of snapshots.
Then the snapshot iterates satisfy
\begin{equation}
\frac{1}{S}\sum_{s=0}^{S-1}\EE\!\left[\bigl\|\nabla f(\tilde w_{s})\bigr\|^2\right]
\;\le\;
\frac{2\bigl(f(w_0)-f^{\star}\bigr)}{\eta S}
\;+\;
\eta L\sigma^2 .
\label{eq:snapshot-rate}
\end{equation}
Moreover, for any $t\in[sk,(s+1)k]$ the interpolated point $\hat w_t$ defined in \eqref{eq:interp} fulfills
\begin{equation}
\EE\!\left[\bigl\|\nabla f(\hat w_t)\bigr\|^2\right]
\;\le\;
2\,\EE\!\left[\bigl\|\nabla f(\tilde w_{s})\bigr\|^2\right]
\;+\;
2\,\EE\!\left[\bigl\|\nabla f(\tilde w_{s+1})\bigr\|^2\right]
\;+\;
L^2k^2\eta^2\sigma^2 .
\label{eq:interp-rate}
\end{equation}
Consequently, choosing $\eta=O(1/\sqrt{T})$ yields
\[
\min_{0\le t\le T}\EE\!\left[\bigl\|\nabla f(\hat w_t)\bigr\|^2\right]=O\!\left(\frac{1}{\sqrt{T}}\right).
\]
\end{theorem}

--------------------------------------------------------------------

\section*{Theoretical Properties}
\begin{itemize}
    \item \textbf{Stability.} Because each update uses a bounded step $\eta\le1/L$, the iterates remain in a region where the $L$‑smoothness inequality holds, guaranteeing that the descent Lemma (Lemma~\ref{lem:descent}) can be applied at every step.
    \item \textbf{Hyper‑parameter sensitivity.}
          \begin{enumerate}[label=(\alph*)]
              \item \emph{Snapshot period $k$.} Larger $k$ reduces memory and communication overhead but inflates the interpolation error term $L^2k^2\eta^2\sigma^2$ in \eqref{eq:interp-rate}. Hence $k$ should be chosen such that $k\eta\sqrt{L\sigma}\ll1$.
              \item \emph{Learning rate $\eta$.} The bound \eqref{eq:snapshot-rate} is minimized for $\eta^{\star}= \sqrt{\frac{2(f(w_0)-f^{\star})}{L\sigma^2 S}}$, reproducing the classic $O(1/\sqrt{T})$ rate.
          \end{enumerate}
    \item \textbf{Comparison to baselines.}
          \begin{itemize}
              \item \emph{Plain SGD.} Standard SGD obtains $\EE\|\nabla f(w_{\text{rnd}})\|^2=O(1/\sqrt{T})$ under the same assumptions. SIS‑SGD achieves the same order while providing a model that can be queried at any intermediate time without recomputing $w_t$.
              \item \emph{Averaged SGD (Polyak‑Ruppert).} Averaging all iterates yields a similar bound but requires storing the whole trajectory. SIS‑SGD stores only $S\approx T/k$ snapshots, offering a favorable memory–accuracy trade‑off.
          \end{itemize}
\end{itemize}

--------------------------------------------------------------------

\section*{Preliminaries (Definitions and Lemmas)}
\begin{lemma}[Descent Lemma]\label{lem:descent}
If $f$ is $L$‑smooth, then for any $x,y\in\RR^{d}$,
\[
f(y)\le f(x)+\inner{\nabla f(x)}{y-x}+\frac{L}{2}\norm{y-x}^2 .
\]
\end{lemma}
\begin{proof}
Standard consequence of $L$‑smoothness; see e.g. \cite{nocedal2006numerical}. \hfill $\square$
\end{proof}

\begin{lemma}[One‑step progress of SGD]\label{lem:sgd-one}
Let $w_{t+1}=w_t-\eta g_t$ with $\EE[g_t\mid w_t]=\nabla f(w_t)$. Then
\begin{equation}
\EE\!\left[f(w_{t+1})\mid w_t\right]
\le f(w_t)-\frac{\eta}{2}\norm{\nabla f(w_t)}^2
      +\frac{\eta L}{2}\sigma^2 .
\label{eq:one-step}
\end{equation}
\end{lemma}
\begin{proof}
Apply Lemma~\ref{lem:descent} with $x=w_t$, $y=w_{t+1}=w_t-\eta g_t$:
\[
f(w_{t+1})\le f(w_t)-\eta\inner{\nabla f(w_t)}{g_t}
               +\frac{L\eta^2}{2}\norm{g_t}^2 .
\tag{Step 1}
\]
Take conditional expectation and use Assumption~\ref{ass:unbiased}:
\[
\EE\!\left[\inner{\nabla f(w_t)}{g_t}\mid w_t\right]
      =\norm{\nabla f(w_t)}^2 .
\tag{Step 2}
\]
Next, decompose $\norm{g_t}^2=\norm{\nabla f(w_t)}^2
        +\norm{g_t-\nabla f(w_t)}^2
        +2\inner{\nabla f(w_t)}{g_t-\nabla f(w_t)}$.
The cross term vanishes after taking expectation (unbiasedness), and
Assumption~\ref{ass:var} bounds the variance term:
\[
\EE\!\left[\norm{g_t}^2\mid w_t\right]
\le \norm{\nabla f(w_t)}^2+\sigma^2 .
\tag{Step 3}
\]
Insert (Step 2) and (Step 3) into (Step 1):
\[
\EE\!\left[f(w_{t+1})\mid w_t\right]
\le f(w_t)-\eta\norm{\nabla f(w_t)}^2
    +\frac{L\eta^2}{2}\bigl(\norm{\nabla f(w_t)}^2+\sigma^2\bigr).
\]
Rearrange terms and use $\eta\le 1/L$ (Assumption~\ref{ass:lr}) to obtain
\[
\EE\!\left[f(w_{t+1})\mid w_t\right]
\le f(w_t)-\frac{\eta}{2}\norm{\nabla f(w_t)}^2
    +\frac{L\eta^2}{2}\sigma^2 .
\]
Since $\eta\le 1/L$, $L\eta^2\le\eta$, which yields the claimed inequality \eqref{eq:one-step}. \hfill $\square$
\end{proof}

--------------------------------------------------------------------

\section*{Main Proof (Step‑by‑Step Derivation)}
We prove Theorem~\ref{thm:main}. The proof is split into two parts: (i) a bound for the snapshot iterates, (ii) a transfer to any interpolated point.

\subsection*{Part (i): Bounding the average squared gradient at snapshots}
\begin{enumerate}
    \item \textbf{Summation of one‑step progress.}  
    Take total expectation in Lemma~\ref{lem:sgd-one} and sum over all $T$ iterations:
    \[
    \EE[f(w_T)]\le f(w_0)-\frac{\eta}{2}\sum_{t=0}^{T-1}\EE\!\bigl[\norm{\nabla f(w_t)}^2\bigr]
                 +\frac{L\eta^2\sigma^2}{2}T .
    \tag{Step 4}
    \]
    Since $f(w_T)\ge f^{\star}$,
    \[
    \frac{\eta}{2}\sum_{t=0}^{T-1}\EE\!\bigl[\norm{\nabla f(w_t)}^2\bigr]
    \le f(w_0)-f^{\star}+ \frac{L\eta^2\sigma^2}{2}T .
    \tag{Step 5}
    \]
    \item \textbf{Extracting snapshot indices.}  
    Recall that snapshots are taken at times $t=sk$ for $s=0,\dots,S-1$.
    Because $k\ge 1$, the set $\{sk\}_{s=0}^{S-1}$ is a subset of $\{0,\dots,T-1\}$,
    therefore
    \[
    \sum_{s=0}^{S-1}\EE\!\bigl[\norm{\nabla f(\tilde w_{s})}^2\bigr]
    \le \sum_{t=0}^{T-1}\EE\!\bigl[\norm{\nabla f(w_t)}^2\bigr].
    \tag{Step 6}
    \]
    \item \textbf{Combining (Step 5) and (Step 6).}
    Dividing by $S$ and using $T\le kS+k$ yields
    \[
    \frac{1}{S}\sum_{s=0}^{S-1}\EE\!\bigl[\norm{\nabla f(\tilde w_{s})}^2\bigr]
    \le \frac{2\bigl(f(w_0)-f^{\star}\bigr)}{\eta S}
          +\eta L\sigma^2\!\left(1+\frac{k}{S}\right) .
    \tag{Step 7}
    \]
    Since $k/S\le 1$ for $T\ge k$, we absorb the factor into the constant and obtain
    \[
    \frac{1}{S}\sum_{s=0}^{S-1}\EE\!\bigl[\norm{\nabla f(\tilde w_{s})}^2\bigr]
    \le \frac{2\bigl(f(w_0)-f^{\star}\bigr)}{\eta S}
          +\eta L\sigma^2 .
    \tag{Step 8}
    \]
    This is exactly the bound \eqref{eq:snapshot-rate}. \hfill $\square$
\end{enumerate}

\subsection*{Part (ii): From snapshots to interpolated points}
Fix $t\in[sk,(s+1)k]$ and write the interpolation \eqref{eq:interp} as
\[
\hat w_t = \tilde w_{s} + \alpha_t(\tilde w_{s+1}-\tilde w_{s}),
\qquad \alpha_t\in[0,1].
\]
Because $\nabla f$ is $L$‑Lipschitz (equivalent to $L$‑smoothness),
\[
\norm{\nabla f(\hat w_t)-\bigl[(1-\alpha_t)\nabla f(\tilde w_{s})+\alpha_t\nabla f(\tilde w_{s+1})\bigr]}
\le L\,\norm{\hat w_t-(1-\alpha_t)\tilde w_{s}-\alpha_t\tilde w_{s+1}}=0 .
\tag{Step 9}
\]
Thus $\nabla f(\hat w_t)$ is exactly the convex combination of the two snapshot gradients:
\[
\nabla f(\hat w_t)=(1-\alpha_t)\nabla f(\tilde w_{s})+\alpha_t\nabla f(\tilde w_{s+1}) .
\tag{Step 10}
\]
Applying Jensen’s inequality to the squared norm of a convex combination,
\[
\bigl\|\nabla f(\hat w_t)\bigr\|^2
\le (1-\alpha_t)\norm{\nabla f(\tilde w_{s})}^2
   +\alpha_t\norm{\nabla f(\tilde w_{s+1})}^2 .
\tag{Step 11}
\]
Taking expectation yields
\[
\EE\!\bigl[\norm{\nabla f(\hat w_t)}^2\bigr]
\le (1-\alpha_t)\EE\!\bigl[\norm{\nabla f(\tilde w_{s})}^2\bigr]
   +\alpha_t\EE\!\bigl[\norm{\nabla f(\tilde w_{s+1})}^2\bigr] .
\tag{Step 12}
\]
To expose the dependence on $k$ and the stochasticity, we bound the difference between consecutive snapshots.
From the SGD update over $k$ steps,
\[
\tilde w_{s+1}-\tilde w_{s}
= -\eta\sum_{j=sk}^{(s+1)k-1} g_j .
\tag{Step 13}
\]
Using the variance bound (Assumption~\ref{ass:var}) and the independence of mini‑batches,
\[
\EE\!\left[\norm{\tilde w_{s+1}-\tilde w_{s}}^2\right]
\le \eta^2 k \sigma^2 .
\tag{Step 14}
\]
Smoothness gives
\[
\norm{\nabla f(\tilde w_{s+1})-\nabla f(\tilde w_{s})}
\le L\,\norm{\tilde w_{s+1}-\tilde w_{s}} .
\tag{Step 15}
\]
Squaring, taking expectation and applying (Step 14) we obtain
\[
\EE\!\bigl[\norm{\nabla f(\tilde w_{s+1})-\nabla f(\tilde w_{s})}^2\bigr]
\le L^2\eta^2 k\sigma^2 .
\tag{Step 16}
\]
Now write
\[
\EE\!\bigl[\norm{\nabla f(\hat w_t)}^2\bigr]
\le 2\EE\!\bigl[\norm{\nabla f(\tilde w_{s})}^2\bigr]
   +2\EE\!\bigl[\norm{\nabla f(\tilde w_{s+1})-\nabla f(\tilde w_{s})}^2\bigr]
\tag{Step 17}
\]
where we used $(a+b)^2\le 2a^2+2b^2$ with $a=\nabla f(\tilde w_{s})$ and
$b=\nabla f(\hat w_t)-\nabla f(\tilde w_{s})$.
Combining (Step 16) with (Step 17) yields exactly \eqref{eq:interp-rate}:
\[
\EE\!\bigl[\norm{\nabla f(\hat w_t)}^2\bigr]
\le 2\EE\!\bigl[\norm{\nabla f(\tilde w_{s})}^2\bigr]
   +2\EE\!\bigl[\norm{\nabla f(\tilde w_{s+1})}^2\bigr]
   +L^2k^2\eta^2\sigma^2 .
\]
\end{enumerate}
\subsection*{Deriving the final rate}
Insert the snapshot bound \eqref{eq:snapshot-rate} into \eqref{eq:interp-rate} and choose a constant step size $\eta = \Theta(1/\sqrt{T})$. Since $S\approx T/k$, the dominant terms become $O(1/(\eta S))=O(1/\sqrt{T})$ and $O(\eta)=O(1/\sqrt{T})$, while the interpolation error $L^2k^2\eta^2\sigma^2 = O(k^2/T)$ is negligible provided $k= o(\sqrt{T})$. Hence
\[
\min_{0\le t\le T}\EE\!\bigl[\norm{\nabla f(\hat w_t)}^2\bigr]=O\!\left(\frac{1}{\sqrt{T}}\right),
\]
which completes the proof of Theorem~\ref{thm:main}. \hfill $\blacksquare$

--------------------------------------------------------------------

\section*{Rate Derivation (With Constants)}
From \eqref{eq:snapshot-rate} we have for any $S\ge1$:
\[
\frac{1}{S}\sum_{s=0}^{S-1}\EE\!\bigl[\norm{\nabla f(\tilde w_{s})}^2\bigr]
\le \underbrace{\frac{2\Delta_f}{\eta S}}_{A}
   +\underbrace{\eta L\sigma^2}_{B},
\quad\text{where }\Delta_f:=f(w_0)-f^{\star}.
\]
Setting $\eta = \sqrt{\frac{2\Delta_f}{L\sigma^2 S}}$ balances $A$ and $B$, yielding
\[
\frac{1}{S}\sum_{s=0}^{S-1}\EE\!\bigl[\norm{\nabla f(\tilde w_{s})}^2\bigr]
\le 2\sqrt{\frac{2L\Delta_f\sigma^2}{S}} .
\]
Since $S\approx T/k$, this is $O\!\bigl(\sqrt{k/T}\bigr)$.  
Plugging into \eqref{eq:interp-rate} and using $k$ fixed gives the final $O(1/\sqrt{T})$ rate for the interpolated iterates.

--------------------------------------------------------------------

\section*{Special Cases}
\begin{enumerate}
    \item \textbf{Deterministic gradients ($\sigma=0$).}  
    The variance term disappears and \eqref{eq:snapshot-rate} reduces to
    $\frac{2\Delta_f}{\eta S}$, i.e. a $O(1/S)=O(k/T)$ rate. In the deterministic smooth setting, SIS‑SGD converges linearly to a stationary point when a diminishing step size $\eta_t\downarrow0$ is used.
    \item \textbf{Infinite snapshot period ($k\to\infty$).}  
    The algorithm collapses to ordinary SGD; the interpolation error term $L^2k^2\eta^2\sigma^2$ blows up, confirming that frequent snapshots are essential for the interpolation guarantee.
    \item \textbf{Convex $f$.}  
    If, in addition, $f$ is convex, the same analysis yields the classic $O(1/T)$ convergence of the function values for the averaged snapshot $\bar w$, matching results for Polyak‑Ruppert averaging.
\end{enumerate}

--------------------------------------------------------------------

\section*{Discussion}
The SIS‑SGD method inherits the $O(1/\sqrt{T})$ stationary‑point guarantee of SGD while providing a lightweight mechanism to obtain a usable model at any intermediate iteration via linear interpolation of sparse snapshots. The theoretical bounds highlight a clear trade‑off: a larger snapshot interval $k$ reduces memory but inflates the interpolation error term $L^2k^2\eta^2\sigma^2$. Choosing $k=O(\sqrt{T})$ balances both effects and retains the optimal rate. Compared with full iterate averaging, SIS‑SGD requires only $O(T/k)$ stored vectors, making it attractive for large‑scale or memory‑constrained training regimes.

\bibliographystyle{plain}
\begin{thebibliography}{9}
\bibitem{nocedal2006numerical}
J.~Nocedal and S.~J. Wright.
\newblock \emph{Numerical Optimization}.
\newblock Springer, 2nd edition, 2006.
\end{thebibliography}

\end{document}